% Ad Familiares 7.3 (ad-familiares-07-03) --- English translation
% Translated by Claude (Anthropic) from the Latin (Perseus).
% Style guide: STYLE.md.

\ciceroLetterOpener{M. CICERO S. D. M. MARIO.}

\ciceroSection{1} Very often, when I am thinking over the common miseries in which we have now spent so many years and, as I see it, will go on spending them, that time we last spent together is apt to come into my mind. Indeed I hold the very day in my memory: for it was on the fourth day before the Ides of May, in the consulship of Lentulus and Marcellus, when I had come to my Pompeian villa late in the day, that you were there to meet me, anxious in mind. What had you anxious was your concern both for my duty and for the danger to me. If I stayed in Italy, you were afraid I should fail in duty; if I set out for the war, my danger weighed on you. At that moment you saw, surely, how I too was so confused that I could not work out what was best to be done. Still, I chose to give way to honour and reputation rather than to take account of my own safety.

\ciceroSection{2} I came to be sorry for that act of mine, not so much because of my own danger as because of the many faults I encountered in the place I had come to: first, no large forces, and none battle-fit; then, apart from the commander and a few besides --- I am speaking of the leading men --- the rest were, in the first place, in the very war itself rapacious, and in the second so cruel of speech that the very victory was a thing of horror; and on top of this the most distinguished men were swimming in the largest debts. What more is there to say? Nothing of any good but the cause. When I had seen this, despairing of victory I began first to urge peace, of which I had always been the advocate; then, when Pompey shrank vehemently from that view, I set about urging him to draw the war out. This he sometimes approved and seemed about to settle in that view, and perhaps would have done so, had he not, after a certain engagement, begun to have confidence in his soldiers. From that moment that great man was no general at all. With raw and pieced-together troops he matched standards against the most seasoned legions; he was beaten in the most disgraceful way; with the very camp lost, he fled alone.

\ciceroSection{3} For myself I made this the end of the war, nor did I think that, having proved unequal at full strength, we should be the stronger when broken. I withdrew from a war in which one had either to fall in line of battle, or be caught in some ambush, or come into the victor's hands, or take refuge with Juba, or take up some place as a kind of exile, or take one's own life by choice. There was, surely, nothing else, if you would not or did not dare to commit yourself to the victor. But of all those troubles I have named, none was more bearable than exile, especially for an innocent man with no shame attached, and especially when one is parted from a city in which one cannot look at anything without grief. As for me, I preferred to be with my own people --- if anyone has anything that is now his own --- and even on what is my own. The things that have happened, I said all of them would happen.

\ciceroSection{4} I came home, not because the conditions of life were the best, but still that, if some shape of a commonwealth survived, I might be as in my country, and if none, as in exile. To take my own life I saw no reason; to wish for it, many reasons. For there is the old saying: when you are not the man you were, there is no reason to want to live. Even so, to be free of fault is a great consolation, especially when I have two things by which I sustain myself, the knowledge of the noblest pursuits and the renown of the greatest deeds; the one of which will never be torn from me while I live, the other not even when I am dead.

\ciceroSection{5} I have written this to you at some length and have been a burden to you, because I have known you to be devoted both to me and to the commonwealth. I wanted my whole policy to be known to you: first, that I never wished any one man to have more power than the commonwealth as a whole; later, when by someone's fault one man was so strong that he could not be resisted, that I wished for peace; with the army lost and the leader on whom the one hope had rested, that I wished to make an end of the war for everyone else as well; and when I could not, that I made an end of it for myself. Now, however, if this is a state, I am a citizen; if it is not, I am an exile, and in no worse a place than if I had betaken myself to Rhodes or Mytilene.

\ciceroSection{6} I had wished to say all this to you face to face; but since the meeting was being put off, I wanted to do it through a letter, so that you should have something to say if you ever came across my detractors. For there are people who, though my destruction would have done the commonwealth no good, take it as a charge against me that I am alive --- and to them, I am very sure, not enough men have died. If they had listened to me, they would still, however unfair the peace, be living honourably; for they would have been beaten in arms, not in cause. There you have a letter perhaps wordier than you would have liked; that I shall take to be how you find it, unless you send back a longer one yourself. If I get done what I want to get done, I shall, as I hope, see you before long.
